% !TEX root = jfla2027_arbres_normatifs.tex

\appendix
\section{Annexes}

\subsection{Figures}

\begin{figure}[htbp]
\begin{Verbatim}[frame=single, label={Code du travail, article R. 1234-2}]
champ d'application MontantIndemnitéDeLicenciement_Loi:
  étiquette plancher_légal_r1234_2
  définition montant_du_plancher
    sous condition années_ancienneté <= 10,0
    conséquence égal à
      montant_du_salaire_de_référence * années_ancienneté * (1,0 / 4,0)

  étiquette plancher_légal_r1234_2
  définition montant_du_plancher
    sous condition années_ancienneté > 10,0
    conséquence égal à
      montant_du_salaire_de_référence * 10,0 * (1,0 / 4,0)
      + montant_du_salaire_de_référence * (années_ancienneté - 10,0)
        * (1,0 / 3,0)
\end{Verbatim}
\begin{Verbatim}[frame=single, label={Convention Syntec (IDCC 1486), article 4.5}]
champ d'application MontantIndemnitéDeLicenciement_CCN_1486
  sous condition
    article_4_5_applicable_VersionApplicableRuptureDuContrat_CCN_1486:
  étiquette barème
  définition montant_du_barème
    sous condition
      faits_rupture.catégorie = Etam
      et années_de_présence <= 10,0
    conséquence égal à
      mois_de_rémunération * années_de_présence * (1,0 / 4,0)

  étiquette barème
  définition montant_du_barème
    sous condition
      faits_rupture.catégorie = IngénieurOuCadre
      et années_de_présence >= 2,0
    conséquence égal à
      mois_de_rémunération * années_de_présence * (1,0 / 3,0)
\end{Verbatim}
\caption{Deux extraits du corpus Catala~\cite{corpuslumo} : le barème légal et
  le barème conventionnel de l'indemnité de licenciement. Commentaires omis et
  mise en page adaptée ; les deux autres cas du barème Syntec (ETAM au-delà de
  dix ans, cadres de moins de deux ans) sont omis.}
\label{fig:catala}
\end{figure}

\subsection{Preuves}

\begin{proof}[\autoref{lem:monotone}]
Munissons \(\Bthree\) de l'ordre d'information \(\leq_i\), dans lequel
\(\unknown\) est au-dessous de \(\true\) et de \(\false\), incomparables entre
eux. Pour un test élémentaire \(t\), ou bien \(\sem{t}_D=\unknown\), ou bien
\(\dep(t)\subseteq\Etablie_D\subseteq\Etablie_{D'}\) et \(\widehat\nu_D\) et
\(\widehat\nu_{D'}\) coïncident sur \(\dep(t)\), d'où \(\sem{t}_{D'}=\sem{t}_D\) ;
dans les deux cas \(\sem{t}_D\leq_i\sem{t}_{D'}\). Les connecteurs forts de
Kleene sont monotones pour \(\leq_i\), de sorte que
\(\sem{\pi}_D\leq_i\sem{\pi}_{D'}\) pour tout prédicat \(\pi\), par induction.
Une valeur décidée étant maximale pour \(\leq_i\), elle est conservée.
\end{proof}

\begin{proof}[\autoref{prop:coh-input}]
Les conditions~1 à~3 et~6 de la \autoref{def:coh-raff} ne dépendent pas du
dossier. Pour la condition~4, une matière ouverte de \(D'\) est ouverte dans
\(D\), ou bien est \(m\) avec \(v=\bot\), abstraite dans \(\rho\) par le
cas~3. Pour la condition~5, une matière indéterminée de \(D'\) est \(m\) avec
\(v=\tree(m)\), feuille de \(\rho\) par le cas~2, pour laquelle la dérivation
\(\tree(m)\to_{D'}^{*}\nu_{D'}(m)\) est vide ; ou bien une matière
\(m''\in\Indet_D\), avec \(\tree(m'')\redD^{*}\nu_D(m'')\). Comme
\(D\enrich D'\), les prédicats décidés dans \(D\) le sont identiquement dans
\(D'\) (\autoref{lem:monotone}), et les frères d'un nœud ne dépendent que de
l'arbre : chaque pas \(\redD\) de cette dérivation est aussi un pas
\(\to_{D'}\).
\end{proof}


\begin{proof}[\autoref{prop:coh-precise}]
Les conditions~4 et~5 de la \autoref{def:coh-raff} passent de \(D\) à
\(D\ominus m\), dont les matières ouvertes et indéterminées sont celles de
\(D\) privées de \(m\), et dont la valuation coïncide avec celle de \(D\) sur
les matières établies. Chaque saisie \(\Input_{\widehat\tree(m')}(\cdot,m')\)
est ensuite valide : \(m'\in\mc M_\rho\) est un enfant de \(m\), donc
distinct de \(\Matter\), indéfini dans le dossier courant, et
\(\widehat\tree(m')\) est son arbre si \(m'\) est concrète, \(\bot\) sinon.
La \autoref{prop:coh-input} conclut.
\end{proof}

\begin{proof}[\autoref{lem:reduits}]
(1) Un pas \((\Cut)\) en \(e\) demande \(\pi(e)(D)=\false\), un pas
\((\Select)\) en \(e=(u,v)\) demande \(\pi(e)(D)=\true\) et
\(\sib_u(v)\neq\emptyset\). D'après le \autoref{lem:monotone}, un prédicat
décidé dans \(D\) a la même valeur dans \(D'\), et les frères d'un nœud ne
dépendent que de l'arbre courant : chaque pas de la dérivation
\(\theta_0\redD^{*}\theta\) est un pas de \(\to_{D'}\).
(2) Les conditions~(a), (d) et~(e) portent sur des couples de nœuds et sont
héritées par tout sous-arbre ; des frères dans \(\theta\) sont frères dans
\(\theta_0\), d'où~(b). Pour l'exhaustivité, soit \(v\) un fils de la racine
de \(\theta_0\) tel que \(\pi((r,v))(D)\neq\false\), qui existe par la
condition~(c) appliquée à \(\theta_0\) et \(D\). Si \(v\) appartient à
\(\theta\), c'est fini. Sinon, \(v\) a été retiré par un pas portant sur une
arête issue de la racine : non par \((\Cut)\) en \((r,v)\), dont le prédicat
n'est pas faux, donc par \((\Select)\) en \((r,v')\) pour un frère \(v'\) tel
que \(\pi((r,v'))(D)=\true\). Ce frère subsiste dans \(\theta\) : \((\Cut)\)
ne s'applique pas à \((r,v')\), et un \((\Select)\) qui le retirerait
exigerait un autre fils de la racine d'arête vraie, ce qu'exclut~(b) pour
\(\theta_0\).
\end{proof}

\begin{proof}[\autoref{prop:norm}]
Terminaison : chaque pas retire au moins un sommet d'un arbre fini. Pour la
confluence locale, on remarque que \(\pi(e)(D)\) ne dépend que de \(D\), et
que tout élagué de \(\nu_D(m)\) dans \(D\) est un élagué de \(\tree(m)\),
qui vérifie donc~(b) d'après le \autoref{lem:reduits}. Soient deux pas en
\(e_i=(u_i,v_i)\), \(i=1,2\), qui retirent les ensembles clos par descendants
\(W_1,W_2\). Si \(v_2\in W_1\), alors \(W_2\subseteq W_1\), et le pas en
\(e_1\) soit s'applique encore dans \(\theta_2\) et donne \(\theta_1\), soit
est un \((\Select)\) dont les sous-arbres frères sont tous inclus dans
\(W_2\), auquel cas \(\theta_1=\theta_2\). Si aucune des deux cibles
n'appartient à l'ensemble de l'autre, les deux pas s'appliquent encore et
commutent. Le cas restant, \(v_1\in W_2\) et \(v_2\in W_1\), force \(v_1\) et
\(v_2\) à être frères avec deux arêtes valant \(\true\), ce qu'exclut~(b). Le
lemme de Newman permet de conclure.
\end{proof}

\begin{proof}[\autoref{prop:coh-prune}]
Les matières établies, ouvertes et indéterminées de \(D'=\Prune(D,m)\) sont
celles de \(D\), aux mêmes valeurs établies, de sorte que \(\to_{D'}\) et
\(\redD\) coïncident. La condition~5 vaut pour \(m\) parce que
\(\tree(m)\redD^{*}\nu_D(m)\redD^{*}\nf{\nu_D(m)}\), et pour les autres
matières indéterminées parce que leur arbre est inchangé.
\end{proof}

\begin{proof}[\autoref{lem:linear}]
(1) Aucune arête de \(\theta'\) ne vaut \(\false\), sinon \((\Cut)\)
s'appliquerait ; toute arête vaut donc \(\true\). L'arbre \(\theta'\) est un
élagué de \(\tree(m)\) dans \(D\), et vérifie donc~(a) et~(b), et
l'exhaustivité de la racine pour \(D\) (\autoref{lem:reduits}). Si un nœud
avait deux fils, les deux arêtes vaudraient \(\true\), ce qui contredirait
l'exclusion mutuelle~(b) ; \(\theta'\) est donc un chemin. La racine conserve
au moins un fils, dont l'arête, non fausse, vaut \(\true\) ; le chemin est
donc non vide. La monotonie des sources découle de~(a).
(2) Par totalité (\autoref{def:applicability-predicate}), une matière de
\(\dep(e)\) n'est pas établie dans \(D\) ; elle est concrète, donc ni ouverte
(condition~4) ni égale à \(\Matter\), et appartient ainsi à
\(\Indet_D\cup\Undef_D\).
\end{proof}

\begin{proof}[\autoref{thm:boucle}]
Pour la phase de concrétisation, associons à chaque matière
abstraite $m$ le poids
\[
w(m)
:=
1+\sum_{\substack{m\refby{\pi}m'\\m'\notin\CMatter}}w(m').
\]
Ces poids sont bien définis puisque le raffinement est fini
et strictement descendant. Posons
\[
W_D:=\sum_{m\in\Ouverte_D}w(m).
\]
Un pas de précision remplace une matière ouverte par certains
de ses enfants : les enfants abstraits sont ouverts, tandis
que les enfants concrets reçoivent directement leur arbre.
Il fait donc strictement décroître $W_D$.
Une saisie factuelle laisse $W_D$ inchangé et fait décroître
le nombre de dépendances de raffinement absentes du dossier.
Le couple
\[
\bigl(W_D,\;|\dep(\rho)\setminus\mc M_D|\bigr)
\]
décroît ainsi lexicographiquement à chaque tour de cette phase.
La totalité des tests garantit qu'un prédicat inconnu possède
une dépendance factuelle manquante. La phase termine donc
avec $\Ouverte_D=\emptyset$.

Pour la phase de résolution, soit $I_D$ l'ensemble croissant
des matières pour lesquelles un arbre du référentiel a déjà
été introduit, et posons
\[
\dep(\mc F_\rho)
:=
\bigcup_{\theta\in\mc F_\rho}\dep(\theta).
\]
Le triplet
\[
\bigl(
|\mathrm{dom}(\mc F_\rho)\setminus I_D|,
\;|\dep(\mc F_\rho)\setminus\mc M_D|,
\;|\Indet_D|
\bigr)
\]
décroît lexicographiquement à chaque tour :
une saisie d'arbre fait décroître sa première composante,
une saisie de valeur sa deuxième, et une résolution sa troisième.
L'élagage termine et est suivi, dans le même tour, d'une saisie
ou d'une résolution.

L'acyclicité des dépendances permet de choisir une matière
indéterminée minimale. Ses dépendances bloquantes sont
indéfinies et peuvent donc être saisies. Lorsqu'elles sont
établies, la totalité des tests, la cohérence du référentiel
et la stabilité des prédicats décidés sous enrichissement
du dossier garantissent un chemin applicable et exécutable ;
la résolution est alors possible. La seconde phase termine donc avec $\Indet_D=\emptyset$,
sans introduire de matière ouverte.

Enfin, les matières de $Q$ ont reçu leurs arbres lors de la
concrétisation et ne sont pas retirées pendant la résolution.
Elles sont donc toutes établies dans le dossier final.
Si $Q$ est vide, la réponse est la famille vide ; sinon,
leurs valeurs constituent la famille annoncée.
\end{proof}
